%%%%%%%%%%%%%%%%%%%%%%%%%%%%%%%%%%%%%%%%%%%%%%%%%%%%%%%%%%%%%%%%%%%%%%
% Lớp 11 | Quan hệ vuông góc trong không gian | Góc giữa đường thẳng và mặt phẳng. Góc nhị diện
% Dạng 1H8?6-7: Toán thực tế về góc đường thẳng, mặt phẳng, góc nhị diện
% Số câu: 1  (sinh tự động bởi tach_cau.py; từ nay thầy sửa trực tiếp file này)
%%%%%%%%%%%%%%%%%%%%%%%%%%%%%%%%%%%%%%%%%%%%%%%%%%%%%%%%%%%%%%%%%%%%%%

%%%%%%%%%%%%%%%% TRẢ LỜI NGẮN %%%%%%%%%%%%%%%%

% nguồn: Ok_TraLoiNgan_Lop11_Lan1
\begin{ex}%[1H8V6-7]%[ma:9182461f]
	Kim tự tháp Memphis tại bang Tennessee (Mỹ) có dạng hình chóp tứ giác đều với chiều cao $98$ m và cạnh đáy $180$ m. Tính số đo góc nhị diện tạo bởi mặt bên và mặt đáy của kim tự tháp đó. (đơn vị đo góc là độ, làm tròn đến hàng phần chục) 
	\shortans{47,4} 
	\loigiai{ 
		\begin{center}
			\begin{tikzpicture}
					\path 
						(0,0)  coordinate[label=left: A] (A)
						(-1,-1) coordinate[label=left: B] (B) 
						(3,0) coordinate[label=right: D]  (D)
						(2,-1) coordinate[label=right: C] (C) 
						(1,3) coordinate[label=left: S]  (S)
						($(A)!0.5!(C)$) coordinate[label=below: O] (O); 
					\draw (D)--(C)--(B); 
					\draw[dashed] (B)--(A)--(D); 
					\draw[dashed]  (A)--(C)
						 (D)--(B);
					 \draw[dashed] (S)--(A);
					\draw  (B)--(S)--(C)
						 (D)--(S);
				\end{tikzpicture}
		\end{center}
		Xét hình chóp tứ giác đều  $S.ABCD$ có chiều cao $98$ m và cạnh đáy $180$ m\\
		Gọi $O$ là tâm  hình vuông $ABCD$ thì $SO\perp (ABCD)\Rightarrow SO\perp AB$. $ \qquad (1)$\\
		Gọi $M$ là trung  điểm $AB$ thì $OM$ là đường trung bình của tam giác $ABC$, \\
		Suy ra $OM=\dfrac{BC}{2}=90$ (m)  và $OM\perp AB.$ $\qquad (2)$\\
		Từ $(1)$ và $(2)$ suy ra $\widehat{SMO}$ là  góc phẳng nhị diện $[S,AB,O]$ \\
		với $$\tan \widehat{SMO}=\dfrac{SO}{OM}=\dfrac{98}{90}=\dfrac{49}{45} \Rightarrow  \widehat{SMO}\approx 47{,}44^{\circ}$$
		Vậy góc nhị diện  tạo bởi mặt bên và mặt đáy của kim tự tháp xấp xỉ $47{,}44^{\circ}$.
	} 
\end{ex}
